\section{Threats to validity and preregistered limitations}
\label{sec:threats}

These are stated before any result exists, so that none of them can be introduced
later as an explanation for an inconvenient outcome.

\begin{itemize}[leftmargin=1.2em]
  \item \textbf{Scope.} Because all 18 databases appear in both partitions, the
  headline estimates unseen-question performance inside a modeled database, not
  cold-start performance on an unmodeled one.
  \item \textbf{Supervision.} The final system uses training failures and
  training gold offline. It is supervised on 231 questions and is not zero-shot.
  \item \textbf{Contamination.} The benchmark is public. The maintainers' design
  claim cannot rule out post-release model exposure, and this project cannot
  inspect training provenance. Paired within-question contrasts are less
  vulnerable than absolute levels.
  \item \textbf{Power.} 101 held-out questions give roughly one percentage point
  of resolution per question, with useful descriptive power but limited power for
  small paired differences and for per-database or rare-condition cells.
  \item \textbf{Metric.} Soft EX deliberately equates some semantically different
  SQL and rejects some extensionally equal results. Accuracy is accuracy under a
  pinned scorer, not proof of semantic equivalence, which is why a second
  comparator is reported.
  \item \textbf{Model parity.} When C4's underlying models cannot be matched,
  C4 minus C3 is confounded and remains a system-level comparison rather than a
  causal estimate of enforcement.
  \item \textbf{Harness.} Published controlled studies place the scaffold effect
  anywhere from 0 to 28 percentage points~\cite{arxiv-2607-22585,arxiv-2606-08529},
  so all condition results are conditional on the disclosed harnesses, and effects
  may appear in error type, cost, retries, or reliability rather than in aggregate
  accuracy.
  \item \textbf{Stochasticity.} Three repetitions characterize some instability
  but do not fully identify provider or temporal variance.
  \item \textbf{Environment parity.} Any mismatch between the database the
  evaluated product queries and the scorer's snapshot can manufacture false gains
  or false failures, which is why fingerprint and canary parity checks precede
  every run.
  \item \textbf{Version drift.} The product, the model providers, and the
  evolving benchmark can all change during a 1,212-trial evaluation. Pinned
  artifacts, timestamps, deterministic interleaving, and the sealed
  generate-then-score boundary mitigate this without eliminating it.
  \item \textbf{Scalability of the modeling claim.} Manual interpretive mappings
  and benchmark-specific exceptions weaken any claim that semantic-model
  construction is automatic, and are counted, not described qualitatively.
  \item \textbf{Composite systems.} A production system that routes across
  multiple opaque model tiers and includes a validation stage cannot be reduced to
  one model label, and will not be reported as though it could.
\end{itemize}
